%% ============================================================
%%  §5 共享契约的谈判成本：mhlo -> StableHLO
%%  要证明的一句话：
%%    把一个方言变成各方都敢依赖的契约，需要版本化、序列化与兼容测试等制度性设施；
%%    这部分成本不体现在 |E| 里，却决定共享层采纳率 p 能否上去。
%%  素材：docs/07-section5-stablehlo.md
%%  建议篇幅：1.5 页
%% ============================================================

\section{共享契约的谈判成本}
\label{sec:contract}

%% 开篇：§3 说明了 MLIR 有能力让层被协商，但没说明"协商成功"需要什么。
%% 本节用一个真实的迁移案例回答。

%% ---- MHLO 的位置问题 ----
%% 它长期作为独立孵化仓库存在于 tensorflow/mlir-hlo，定位是"独立于 TF 与 XLA 的端到端流程"、
%% "把技术上游化的孵化器"。把它当框架与编译器之间的接口，等于把跨组织接口放在一个项目的内部 IR 上。
%% 现状：官方已标 deprecated，建议迁移（措辞是 deprecated + 迁移建议，不是"设计不好"）。

%% ---- StableHLO 的增量：全是契约设施，不是算子语义 ----
%% 官方自述："基于 MHLO，并增加了序列化与版本化"，以此提供前后向兼容保证。
\begin{table}[t]
  \caption{StableHLO 相对 MHLO 增加的契约设施。新增项均为"制度性"机制，
  而非算子语义改造。}
  \label{tab:contract}
  \small
  \begin{tabularx}{\linewidth}{@{}l X@{}}
    \toprule
    机制 & 具体做法 \\
    \midrule
    定位 & 明确充当 ML 框架与编译器之间的 portability layer \\
    序列化 & portable artifact 用 MLIR bytecode；CLI / C++ / Python 三种入口 \\
    版本号 & \texttt{\#.\#.\#}：minor 随 opset 或序列化格式变更递增；patch 随发布递增，每周一次 \\
    不可回填 & 版本被 supersede 后不能再添加内容，变更必须开新 minor（旧版本语义被冻结） \\
    VHLO & 把"版本化"本身做成一个方言；序列化时 StableHLO $\rightarrow$ 目标版本的 VHLO \\
    兼容窗口 & \textbf{5 年向后兼容} + \textbf{2 年向前兼容} \\
    CI 门槛 & 每个 PR 都测双向兼容；目标版本取至少 4 / 12 周前被 supersede 的版本 \\
    显式排除 & prettyprinted IR、未规格化特性、bug compatibility 均\textbf{不}保证 \\
    \bottomrule
  \end{tabularx}
\end{table}

%% ---- 三条推论 ----
%% A. 契约的成本是治理成本：上表没有一行是"再实现一个算子"；成本在 RFC、每周发版、
%%    一个 VHLO 方言、逐 PR 的兼容测试套。
%% B. 这套治理能力不在 MLIR 里：compatibility.md 的 Future work 明确说要"把兼容性测试套
%%    贡献回 MLIR 上游"（issue #1632）。在贡献之前，契约是 StableHLO 自己立法、自己执行、自己承担。
%%    => 这正是 §4 结论"收益条件性"的机制解释：MLIR 提供机制，契约由采用者建立。
%% C. 契约必须显式排除才能成立：out-of-scope 三条说明"边界可声明，但不等于边界免费"，
%%    接得上 §3 的机制 ④。

%% 收束句（建议原话）：
%% 由此可见，MLIR 让层可被协商，但把某一层变成各方都敢依赖的契约，需要采用者自建版本、
%% 序列化与兼容测试制度；这部分成本不体现在转换数量 |E| 里，却决定了 p 能否真正上去。
